\documentclass[10pt,a4paper]{amsart}
\usepackage[margin=19mm]{geometry}
\usepackage{amsmath,amssymb,mathtools}
\usepackage[scr=rsfs,cal=euler]{mathalfa}
\usepackage{xfrac}
\usepackage[unicode,hidelinks,bookmarks=false]{hyperref}
\newcommand{\ie}{i.e.\ }
\newcommand{\cf}{cf.\ }
\newcommand{\resp}{respectively\ }
\newcommand{\Subsection}{\subsection}
\providecommand{\nobreakdash}{\nobreak\hbox{-}}
\setlength{\emergencystretch}{2em}
\multlinegap0pt
\usepackage{bm}
\usepackage{bbm}
\usepackage{dsfont}
\usepackage{xspace}
\usepackage{cancel}
\usepackage{mathtools}
\usepackage{amsthm}
\newtheorem{theorem}{Theorem}
\title[New Generalized Dual Banach Frames]{New Generalized Dual Banach Frames (Theorem 3.4)}
\author{S. Shahsavand, M. Abolghasemi}
\date{Source: arXiv:2609.28803v1 (CC BY 4.0)}
\begin{document}
\maketitle

\noindent\emph{Source and licence.} New Generalized Dual Banach Frames. Version arXiv:2609.28803v1 (CC BY 4.0), distributed under the Creative Commons Attribution 4.0 International licence, as recorded in the arXiv licence metadata for this exact version and shown on the abstract page https://arxiv.org/abs/2609.28803v1. Full rights notice: licenses/arxiv-2609.28803-CC-BY-4.0.txt.\par\smallskip

\noindent\textit{Editorial scope.} Counted unit: the Theorem printed as Theorem 3.4 of the article (source0.tex lines 1606--1649, source label \texttt{None}), delivered together with the article's own complete original proof of it (source0.tex lines 1650--1799, one \texttt{proof} environment of 150 lines). No mathematics is written, repaired or re-derived: statement and proof are the article's own text line for line. Required context retained uncounted: none. Every definition, display and label of those lines is retained unchanged. Omitted from this package and not redistributed: 1--1605, 1800--2340. Nothing inside a retained excerpt is omitted or altered: every retained body is delivered exactly as the article's lines are written, so no omission receipt of any kind accompanies this package. Citations: the retained excerpts carry no citation command at all, so no bibliography entry of the article is transcribed and no reference is redistributed. Reference adaptation: none was needed and none was made; every reference inside the delivered unit resolves inside it, so no unresolved reference or citation is produced. Printed slips kept literal: no printed word, symbol, hypothesis or number of the counted statement or of its proof was altered. Layout and class: the article's own class is \texttt{article} with packages including inputenc, amsmath, amssymb, amsthm, geometry, mathrsfs, hyperref, enumitem, setspace, natbib; the wrapper is the lane's pinned \texttt{amsart} editorial wrapper at 10pt on A4, which restates the theorem-like environments the retained text needs and adds the article's own macro definitions that the retained text needs (none), with extra wrapper packages (bm, bbm, dsfont, xspace, cancel, mathtools). The delivered numbering is the unit's own. Assets: no retained span contains a figure environment, a graphics-inclusion command, a PDF-page inclusion, a table or a TikZ picture; no image file is copied into this package, the package declares no asset dependency and no image file is redistributed with it. Licence: version arXiv:2609.28803v1 of this article is distributed under the Creative Commons Attribution 4.0 International licence, as recorded in the arXiv licence metadata for this exact version and shown on the abstract page https://arxiv.org/abs/2609.28803v1. A full rights notice is delivered at licenses/arxiv-2609.28803-CC-BY-4.0.txt. Characters: every retained excerpt is pure ASCII, so the delivered engine, XeLaTeX with its default Latin Modern text font, reports no missing character.
\begin{theorem}
	Assume that the sequences \(\{f_i\}_{i\in I}\subset X^*\) and
	\(\{x_i\}_{i\in I}\subset X\) satisfy the \(X_a\)-Bessel condition in
	\(X^*\) and \(X\), respectively. Let
	\[
	U:X\to X_a,\qquad V:X_a\to X
	\]
	denote the corresponding analysis and synthesis operators. Then the
	following statements hold:
	
	\begin{enumerate}
		\item If \(\{f_i\}_{i\in I}\) is a Banach dual frame for
		\(\{x_i\}_{i\in I}\), then \(\{f_i\}_{i\in I}\) is an approximately
		dual Banach frame for \(\{x_i\}_{i\in I}\).
		
		\item If \(\{f_i\}_{i\in I}\) is an approximately dual Banach frame
		for \(\{x_i\}_{i\in I}\), then \(\{f_i\}_{i\in I}\) is a pseudo-dual
		Banach frame for \(\{x_i\}_{i\in I}\).
		
		\item Suppose that \(\{f_i\}_{i\in I}\subset X^*\) is a
		pseudo-dual Banach frame for \(\{x_i\}_{i\in I}\subset X\), and let
		\(G\in B(X)\) be invertible. Then \(\{G^*{}^{-1}f_i\}_{i\in I}\)
		is a pseudo-dual Banach frame for \(\{Gx_i\}_{i\in I}\).
		
		\item If \(\{f_i\}_{i\in I}\) is a pseudo-dual Banach frame for
		\(\{x_i\}_{i\in I}\), then it is a \(g\)-dual Banach frame for
		\(\{x_i\}_{i\in I}\), associated with the operator
		\[
		A=(VU)^{-1}.
		\]
		
		\item If \(\{f_i\}_{i\in I}\) is a pseudo-dual Banach frame for
		\(\{x_i\}_{i\in I}\), then \(\{f_i\}_{i\in I}\) is a Banach dual
		frame for the sequence
		\[
		\{(VU)^{-1}x_i\}_{i\in I}.
		\]
		
		\item The sequence \(\{x_i\}_{i\in I}\) is a \(g\)-dual Banach frame
		for \(\{f_i\}_{i\in I}\) if and only if
		\(\{x_i\}_{i\in I}\) is a pseudo-dual Banach frame for
		\(\{f_i\}_{i\in I}\).
	\end{enumerate}
\end{theorem}

\begin{proof}
	Let
	\[
	U:X\to X_a,\qquad V:X_a\to X
	\]
	be the analysis and synthesis operators associated with
	\(\{f_i\}_{i\in I}\) and \(\{x_i\}_{i\in I}\), respectively.
	
	\begin{enumerate}
		\item Suppose that \(\{f_i\}_{i\in I}\) is a Banach dual frame for
		\(\{x_i\}_{i\in I}\). By the definition of Banach duality, we have
		\[
		VU=I_X.
		\]
		Therefore,
		\[
		\|I_X-VU\|=0<1.
		\]
		Hence, \(\{f_i\}_{i\in I}\) is an approximately dual Banach frame
		for \(\{x_i\}_{i\in I}\).
		
		\item Assume that \(\{f_i\}_{i\in I}\) is an approximately dual
		Banach frame for \(\{x_i\}_{i\in I}\). Then
		\[
		\|I_X-VU\|<1.
		\]
		By the Neumann series theorem, \(VU\) is invertible on \(X\), with
		\[
		(VU)^{-1}
		=
		\sum_{n=0}^{\infty}(I_X-VU)^n.
		\]
		Thus, \(\{f_i\}_{i\in I}\) is a pseudo-dual Banach frame for
		\(\{x_i\}_{i\in I}\).
		
		\item Suppose that \(\{f_i\}_{i\in I}\) is a pseudo-dual Banach
		frame for \(\{x_i\}_{i\in I}\), and let \(G\in B(X)\) be invertible.
		Define
		\[
		\widetilde{x}_i=Gx_i,
		\qquad
		\widetilde{f}_i=(G^{-1})^*f_i.
		\]
		The corresponding analysis operator is
		\[
		\widetilde{U}x
		=
		\bigl(\widetilde{f}_i(x)\bigr)_{i\in I}
		=
		\bigl(f_i(G^{-1}x)\bigr)_{i\in I}
		=
		U(G^{-1}x),
		\]
		while the synthesis operator satisfies
		\[
		\widetilde{V}(c_i)_{i\in I}
		=
		\sum_{i\in I}c_iGx_i
		=
		G V(c_i)_{i\in I}.
		\]
		Consequently,
		\[
		\widetilde{V}\widetilde{U}
		=
		GV G^{-1}.
		\]
		Since \(VU\) is invertible and \(G\) is invertible,
		\(GVG^{-1}\) is invertible. Hence
		\(\{\widetilde{f}_i\}_{i\in I}\) is a pseudo-dual Banach frame
		for \(\{\widetilde{x}_i\}_{i\in I}\).
		
		\item Suppose that \(\{f_i\}_{i\in I}\) is a pseudo-dual Banach
		frame for \(\{x_i\}_{i\in I}\). Then \(VU\) is invertible. Put
		\[
		A=(VU)^{-1}.
		\]
		For every \(x\in X\),
		\[
		\sum_{i\in I}f_i(Ax)x_i
		=
		VU(Ax)
		=
		VU(VU)^{-1}x
		=
		x.
		\]
		Therefore, \(\{f_i\}_{i\in I}\) is a \(g\)-dual Banach frame for
		\(\{x_i\}_{i\in I}\), associated with the invertible operator
		\(A=(VU)^{-1}\).
		
		\item Let \(\{f_i\}_{i\in I}\) be a pseudo-dual Banach frame for
		\(\{x_i\}_{i\in I}\). Define
		\[
		\widetilde{x}_i=(VU)^{-1}x_i.
		\]
		Then, for every \(x\in X\),
		\[
		\begin{aligned}
			\sum_{i\in I}f_i(x)\widetilde{x}_i
			&=
			\sum_{i\in I}f_i(x)(VU)^{-1}x_i\\
			&=
			(VU)^{-1}
			\sum_{i\in I}f_i(x)x_i\\
			&=
			(VU)^{-1}VUx\\
			&=x.
		\end{aligned}
		\]
		Hence, \(\{f_i\}_{i\in I}\) is a Banach dual frame for
		\(\{\widetilde{x}_i\}_{i\in I}\).
		
		\item Suppose that \(\{x_i\}_{i\in I}\) is a \(g\)-dual Banach frame
		for \(\{f_i\}_{i\in I}\). Then there exists an invertible operator
		\(A\in B(X^*)\) such that
		\[
		f
		=
		\sum_{i\in I}f_i(Af)x_i,
		\qquad f\in X^*.
		\]
		Thus, the corresponding composition of the analysis and synthesis
		operators is invertible on \(X^*\), and consequently
		\(\{x_i\}_{i\in I}\) is a pseudo-dual Banach frame for
		\(\{f_i\}_{i\in I}\).
		
		Conversely, suppose that \(\{x_i\}_{i\in I}\) is a pseudo-dual
		Banach frame for \(\{f_i\}_{i\in I}\). Then the corresponding
		operator \(U^*V^*\) is invertible on \(X^*\). Setting
		\[
		A=(U^*V^*)^{-1},
		\]
		we obtain
		\[
		U^*V^*Af=f,
		\qquad f\in X^*.
		\]
		Hence,
		\[
		f
		=
		\sum_{i\in I}f_i(Af)x_i,
		\qquad f\in X^*,
		\]
		which is precisely the \(g\)-duality relation. Therefore,
		\(\{x_i\}_{i\in I}\) is a \(g\)-dual Banach frame for
		\(\{f_i\}_{i\in I}\).
	\end{enumerate}
\end{proof}

\end{document}
